%!TEX program = lualatex
\documentclass[12pt,a4paper]{article}

% =========================================================
% CONFIGURAÇÃO BÁSICA
% =========================================================
\usepackage[brazil]{babel}
\usepackage{fontspec}
\setmainfont{Latin Modern Roman}

\usepackage[a4paper,margin=2.5cm]{geometry}
\usepackage{microtype}
\usepackage{setspace}
\usepackage{parskip}
\usepackage{enumitem}
\usepackage{xcolor}
\usepackage{hyperref}
\usepackage{csquotes}
\usepackage{fancyvrb}

% =========================================================
% BIBLIOGRAFIA
% =========================================================
\usepackage[
    backend=biber,
    style=authoryear,
    sorting=nyt
]{biblatex}

\addbibresource{referencias.bib}

% =========================================================
% COMANDOS AUXILIARES
% =========================================================
\newcommand{\arquivo}[1]{\texttt{#1}}
\newcommand{\comando}[1]{\texttt{#1}}

\DefineVerbatimEnvironment{bloco}{Verbatim}{
    fontsize=\small,
    frame=single,
    rulecolor=\color{black},
    numbers=none
}

\title{Guia Prático para usar \texttt{biblatex} + \texttt{biber} no LuaLaTeX}
\author{}
\date{\today}

\begin{document}

\maketitle
\onehalfspacing

\section{Objetivo}

Este documento registra o procedimento correto para usar referências bibliográficas em um projeto \LaTeX\ com:

\begin{itemize}[leftmargin=1.5cm]
    \item \textbf{LuaLaTeX} como motor de compilação;
    \item \textbf{biblatex} como pacote de bibliografia;
    \item \textbf{biber} como processador do arquivo \arquivo{.bib}.
\end{itemize}

A lógica de funcionamento é:

\begin{enumerate}[leftmargin=1.5cm]
    \item criar um arquivo de referências, por exemplo \arquivo{referencias.bib};
    \item inserir no texto comandos como \comando{\textbackslash cite\{...\}};
    \item compilar com \comando{lualatex};
    \item processar a bibliografia com \comando{biber};
    \item recompilar com \comando{lualatex}.
\end{enumerate}

\section{Estrutura mínima do projeto}

Uma estrutura simples e funcional é:

\begin{bloco}
Geo_iii/
|-- notas_sobre_geo_iii.tex
|-- referencias.bib
\end{bloco}

O arquivo principal do texto é \arquivo{notas_sobre_geo_iii.tex}.  
O banco de referências é \arquivo{referencias.bib}.

\section{Modelo correto do preâmbulo}

No preâmbulo do arquivo principal, o trecho essencial é:

\begin{bloco}
\usepackage[
    backend=biber,
    style=authoryear,
    sorting=nyt
]{biblatex}

\addbibresource{referencias.bib}
\end{bloco}

Esse trecho informa ao \LaTeX\ que:

\begin{itemize}[leftmargin=1.5cm]
    \item o pacote de bibliografia usado será o \texttt{biblatex};
    \item o processador será o \texttt{biber};
    \item as referências estão armazenadas em \arquivo{referencias.bib}.
\end{itemize}

\section{Como citar no texto}

Supondo que existam no arquivo \arquivo{referencias.bib} as chaves

\begin{bloco}
zerfass2008foraminiferos
milhomem2023baciassedimentares
\end{bloco}

as citações no texto podem ser feitas assim:

\begin{bloco}
Segundo \cite{zerfass2008foraminiferos}, os foraminíferos ...

Uma introdução ao tema pode ser vista em \cite{milhomem2023baciassedimentares}.
\end{bloco}

Também é possível citar mais de uma referência ao mesmo tempo:

\begin{bloco}
\cite{zerfass2008foraminiferos,milhomem2023baciassedimentares}
\end{bloco}

\section{Como imprimir a bibliografia}

No final do documento, antes de \comando{\textbackslash end\{document\}}, deve aparecer:

\begin{bloco}
\printbibliography
\end{bloco}

Sem esse comando, a lista de referências não será exibida no PDF.

\section{O que não deve ser misturado}

Se você usa \texttt{biblatex}, \textbf{não deve misturar} com o esquema antigo do BibTeX, isto é, não deve usar comandos como:

\begin{bloco}
\bibliographystyle{plain}
\bibliography{referencias}
\end{bloco}

Esses comandos pertencem ao fluxo antigo de BibTeX.  
No caso de \texttt{biblatex}, o correto é:

\begin{bloco}
\usepackage[backend=biber]{biblatex}
\addbibresource{referencias.bib}
...
\printbibliography
\end{bloco}

\section{Sequência correta de compilação}

No PowerShell, dentro da pasta do projeto, o fluxo recomendado é:

\begin{bloco}
cd D:\LaTex_Projects\Geo_iii
del *.aux, *.bcf, *.bbl, *.blg, *.run.xml -ErrorAction SilentlyContinue
lualatex notas_sobre_geo_iii.tex
biber notas_sobre_geo_iii
lualatex notas_sobre_geo_iii.tex
lualatex notas_sobre_geo_iii.tex
\end{bloco}

\subsection*{Interpretação da sequência}

\begin{enumerate}[leftmargin=1.5cm]
    \item o primeiro \comando{lualatex} gera os arquivos auxiliares;
    \item o \comando{biber} lê o arquivo \arquivo{.bcf} e produz a bibliografia;
    \item os dois \comando{lualatex} finais resolvem as citações e atualizam a lista bibliográfica.
\end{enumerate}

\section{Como usar no terminal do VS Code}

No terminal integrado do VS Code, você deve digitar \textbf{somente o comando}, e nunca o prompt.

Exemplo de prompt:

\begin{bloco}
PS D:\LaTex_Projects\Geo_iii>
\end{bloco}

Isso \textbf{não} deve ser digitado.  
O que deve ser digitado é apenas:

\begin{bloco}
cd D:\LaTex_Projects\Geo_iii
\end{bloco}

ou, se já estiver na pasta correta:

\begin{bloco}
biber notas_sobre_geo_iii
\end{bloco}

\section{Erros mais comuns}

\subsection*{1. Citation undefined}

Mensagem típica:

\begin{bloco}
LaTeX: Citation 'zerfass2008foraminiferos' undefined
\end{bloco}

Significa que a chave da referência ainda não foi resolvida.  
As causas mais comuns são:

\begin{itemize}[leftmargin=1.5cm]
    \item a chave não existe no arquivo \arquivo{.bib};
    \item o \comando{biber} não foi executado;
    \item o documento foi recompilado em ordem errada;
    \item houve falha em algum arquivo auxiliar.
\end{itemize}

\subsection*{2. Please (re)run Biber}

Mensagem típica:

\begin{bloco}
Package biblatex: Please (re)run Biber on the file:
notas_sobre_geo_iii
and rerun LaTeX afterwards.
\end{bloco}

Essa mensagem deve ser interpretada literalmente: é preciso rodar

\begin{bloco}
biber notas_sobre_geo_iii
\end{bloco}

e depois recompilar com \comando{lualatex}.

\subsection*{3. Empty bibliography}

Mensagem típica:

\begin{bloco}
LaTeX: Empty bibliography.
\end{bloco}

Isso geralmente significa uma das seguintes coisas:

\begin{itemize}[leftmargin=1.5cm]
    \item o arquivo \arquivo{referencias.bib} não foi encontrado;
    \item o \comando{biber} não rodou corretamente;
    \item o documento não conseguiu montar a bibliografia a partir das entradas disponíveis.
\end{itemize}

\section{Checklist de diagnóstico}

Quando a bibliografia falhar, verificar na seguinte ordem:

\begin{enumerate}[leftmargin=1.5cm]
    \item o arquivo \arquivo{referencias.bib} existe na pasta do projeto;
    \item o nome dado em \comando{\textbackslash addbibresource\{referencias.bib\}} está correto;
    \item as chaves usadas em \comando{\textbackslash cite\{...\}} existem exatamente no \arquivo{.bib};
    \item o comando \comando{biber notas_sobre_geo_iii} foi executado;
    \item o documento foi recompilado com \comando{lualatex} depois do \comando{biber};
    \item não há mistura de \texttt{biblatex} com comandos do BibTeX antigo.
\end{enumerate}

\section{Modelo mínimo funcional do arquivo principal}

A seguir está um exemplo mínimo de arquivo principal funcional:

\begin{bloco}
%!TEX program = lualatex
\documentclass{article}
\usepackage[brazil]{babel}
\usepackage{fontspec}
\usepackage[backend=biber,style=authoryear]{biblatex}
\addbibresource{referencias.bib}

\begin{document}

Texto de teste com citação \cite{zerfass2008foraminiferos}.

\printbibliography

\end{document}
\end{bloco}

\section{Modelo mínimo funcional do arquivo \arquivo{referencias.bib}}

\begin{bloco}
@online{milhomem2023baciassedimentares,
  author  = {Milhomem, Autor Exemplo},
  title   = {Bacias sedimentares: introdução e conceitos fundamentais},
  year    = {2023},
  url     = {https://exemplo.org/bacias-sedimentares},
  urldate = {2026-03-17}
}

@article{zerfass2008foraminiferos,
  author       = {Zerfass, Autor Exemplo},
  title        = {Foraminíferos e sua aplicação em estudos geológicos},
  journaltitle = {Revista Exemplo de Geociências},
  year         = {2008},
  volume       = {12},
  number       = {3},
  pages        = {45--60}
}
\end{bloco}

\section{Observação sobre codificação}

Como o projeto está sendo compilado com LuaLaTeX, os arquivos devem ser mantidos em:

\begin{itemize}[leftmargin=1.5cm]
    \item UTF-8;
    \item sem BOM.
\end{itemize}

Nesse caso, não é necessário usar \comando{\textbackslash usepackage[utf8]\{inputenc\}}.

\section{Procedimento recomendado para o projeto real}

Para o arquivo \arquivo{notas_sobre_geo_iii.tex}, o procedimento recomendado é:

\begin{enumerate}[leftmargin=1.5cm]
    \item manter \arquivo{referencias.bib} na mesma pasta do arquivo principal;
    \item usar \comando{\textbackslash addbibresource\{referencias.bib\}};
    \item citar com \comando{\textbackslash cite\{...\}};
    \item incluir \comando{\textbackslash printbibliography} no fim do documento;
    \item compilar pela sequência \comando{lualatex} $\rightarrow$ \comando{biber} $\rightarrow$ \comando{lualatex} $\rightarrow$ \comando{lualatex}.
\end{enumerate}

\section{Conclusão}

O problema mais frequente não está na chave da citação, mas no fluxo de compilação.

Se:

\begin{itemize}[leftmargin=1.5cm]
    \item o arquivo \arquivo{.bib} existe,
    \item a chave existe,
    \item o preâmbulo está correto,
\end{itemize}

então a correção quase sempre consiste em rodar corretamente:

\begin{bloco}
lualatex arquivo.tex
biber arquivo
lualatex arquivo.tex
lualatex arquivo.tex
\end{bloco}

\printbibliography

\end{document}